\chapter{高等数学：二重积分}

\section{Riemann 和、对称性与区域重心}

\begin{problemblock}
\textbf{例1.} 已知
\[
\lim_{n\to\infty}\sum_{i=1}^n\sum_{j=1}^n
\frac{4}{(2n+2i-1)(2n+2j-1)}=k
\]
（\(k\) 为常数），则下列说法中错误的是（\quad）

A. \(\displaystyle \lim_{n\to\infty}\sum_{i=1}^n\sum_{j=1}^n
\frac1{(n+i)(n+j)}=k\)

B. \(\displaystyle \lim_{n\to\infty}\sum_{i=1}^{2n}\sum_{j=1}^{2n}
\frac1{(2n+i)(2n+j)}=k\)

C. \(\displaystyle \lim_{n\to\infty}\sum_{i=1}^n\sum_{j=1}^{2026n}
\frac1{(n+i)(2026n+j)}=k\)

D. \(\displaystyle \lim_{n\to\infty}\sum_{i=1}^{2026n}\sum_{j=1}^n
\frac1{(n+2026i)(n+j)}=\frac{k}{2026}\)
\end{problemblock}
\begin{solutionblock}
\analysis{这类双重和本质是两个一维和的乘积。先识别每个因子对应的对数极限。}
原式可写成
\[
\left[\sum_{i=1}^n\frac{2}{2n+2i-1}\right]^2.
\]
它对应
\[
\int_0^1\frac{dx}{1+x}=\ln2,
\]
所以
\[
k=(\ln2)^2.
\]
A 中
\[
\sum_{i=1}^n\frac1{n+i}\to\ln2,
\]
故极限为 \(k\)，A 正确。

B 中
\[
\sum_{i=1}^{2n}\frac1{2n+i}\to\ln2,
\]
故极限也为 \(k\)，B 正确。

C 中
\[
\sum_{i=1}^n\frac1{n+i}\to\ln2,\qquad
\sum_{j=1}^{2026n}\frac1{2026n+j}\to\ln2,
\]
故极限仍为 \(k\)，C 正确。

D 中第一因子为
\[
\sum_{i=1}^{2026n}\frac1{n+2026i}
\to
\int_0^{2026}\frac{dt}{1+2026t}
=\frac1{2026}\ln(1+2026^2),
\]
并非 \(\frac1{2026}\ln2\)。因此 D 错误。
\answer{D}
\examnote{双重和若可分离，先拆成两个一维 Riemann 和；注意分母中 \(2026i\) 与 \(2026n+i\) 的差别。}
\end{solutionblock}

\begin{problemblock}
\textbf{例2.} 设函数 \(f(x,y)\) 在区域
\[
D=\{(x,y)\mid 0\le x\le1,\ 0\le y\le1\}
\]
上连续，且 \(f(x,y)=f(y,x)\)，则
\[
\iint_D f(x,y)\,dxdy=(\quad)
\]
\begin{enumerate}[label=\Alph*., leftmargin=2.2em]
\item \(\displaystyle 2\lim_{n\to\infty}\sum_{i=1}^n\sum_{j=n+1-i}^n
f\left(\frac in,\frac jn\right)\frac1{n^2}\)
\item \(\displaystyle \frac12\lim_{n\to\infty}\sum_{i=1}^n\sum_{j=1}^n
f\left(\frac in,\frac jn\right)\frac1{n^2}\)
\item \(\displaystyle 2\lim_{n\to\infty}\sum_{i=1}^{2n}\sum_{j=1}^{2n+1-i}
f\left(\frac i{2n},\frac j{2n}\right)\frac1{n^2}\)
\item \(\displaystyle \frac12\lim_{n\to\infty}\sum_{i=1}^{2n}\sum_{j=1}^{i}
f\left(\frac i{2n},\frac j{2n}\right)\frac1{n^2}\)
\end{enumerate}
\end{problemblock}
\begin{solutionblock}
\analysis{条件 \(f(x,y)=f(y,x)\) 说明正方形关于直线 \(y=x\) 的两个三角形积分相等。}
在 \(2n\times2n\) 网格下，单个小矩形面积为
\[
\frac1{(2n)^2}=\frac1{4n^2}.
\]
选项 D 中的和覆盖三角形
\[
0\le y\le x\le1,
\]
但权重是 \(1/n^2=4\cdot 1/(2n)^2\)。因此
\[
\lim_{n\to\infty}\sum_{i=1}^{2n}\sum_{j=1}^{i}
f\left(\frac i{2n},\frac j{2n}\right)\frac1{n^2}
=4\iint_{0\le y\le x\le1}f(x,y)\,dxdy.
\]
再乘以 \(\frac12\)，得
\[
2\iint_{0\le y\le x\le1}f(x,y)\,dxdy.
\]
由于 \(f(x,y)=f(y,x)\)，该值正好等于整个正方形上的积分。
\answer{D}
\examnote{对称函数在正方形上下两个三角形上的积分相等，但网格权重必须同步检查。}
\end{solutionblock}

\begin{problemblock}
\textbf{例3.} 设 \(D\) 是 \(xOy\) 平面上以 \((1,1),(-1,1),(-1,-1)\) 为顶点的三角形区域，\(D_1\) 是 \(D\) 在第一象限的部分，则
\[
\iint_D\bigl(xy+\cos x\sin y\bigr)\,dxdy=(\quad)
\]

A. \(\displaystyle 2\iint_{D_1}\cos x\sin y\,dxdy\)
\quad
B. \(\displaystyle 2\iint_{D_1}xy\,dxdy\)

C. \(\displaystyle 4\iint_{D_1}(xy+\cos x\sin y)\,dxdy\)
\quad
D. \(0\)
\end{problemblock}
\begin{solutionblock}
\analysis{区域 \(D\) 可表示为 \(-1\le x\le1,\ x\le y\le1\)。直接利用奇偶性和第一象限部分比较。}
先看 \(xy\) 项：
\[
\iint_D xy\,dxdy
=\int_{-1}^1x\left(\int_x^1y\,dy\right)dx
=\int_{-1}^1\frac{x(1-x^2)}2\,dx=0.
\]
再看
\[
\iint_D\cos x\sin y\,dxdy
=\int_{-1}^1\cos x\left(\int_x^1\sin y\,dy\right)dx.
\]
而 \(D_1=\{(x,y)\mid 0\le x\le y\le1\}\)。同样计算可知
\[
\iint_D\cos x\sin y\,dxdy
=2\iint_{D_1}\cos x\sin y\,dxdy.
\]
故原积分等于选项 A。
\answer{A}
\examnote{区域不一定整体关于坐标轴对称，必要时直接写出积分限验证奇偶项是否消失。}
\end{solutionblock}

\begin{problemblock}
\textbf{例4.} 设
\[
D=\{(x,y)\mid 2(x-1)^2+3(y-2)^2\le6\},
\]
求
\[
\iint_D(x+y)\,dxdy.
\]
\end{problemblock}
\begin{solutionblock}
\analysis{椭圆关于中心 \((1,2)\) 对称，线性函数在对称区域上的积分等于中心函数值乘面积。}
区域可写为
\[
\frac{(x-1)^2}{3}+\frac{(y-2)^2}{2}\le1.
\]
椭圆中心为 \((1,2)\)，面积为
\[
S=\pi\sqrt3\sqrt2=\pi\sqrt6.
\]
由中心对称性，
\[
\iint_Dx\,dxdy=1\cdot S,\qquad
\iint_Dy\,dxdy=2\cdot S.
\]
所以
\[
\iint_D(x+y)\,dxdy=3S=3\pi\sqrt6.
\]
\answer{\(3\pi\sqrt6\)}
\examnote{对称区域上线性函数积分：\(\iint_D(ax+by+c)= [a x_c+b y_c+c]\,S_D\)。}
\end{solutionblock}

\begin{problemblock}
\textbf{例5.} 计算
\[
\iint_D(x+y)\,dxdy,
\qquad
D=\{(x,y)\mid x^2+y^2\le x+y+1\}.
\]
\end{problemblock}
\begin{solutionblock}
\analysis{配方后得到圆盘，仍用中心乘面积。}
配方：
\[
x^2-x+y^2-y\le1
\]
即
\[
\left(x-\frac12\right)^2+\left(y-\frac12\right)^2\le\frac32.
\]
圆心为 \((1/2,1/2)\)，半径为 \(\sqrt{3/2}\)，面积为
\[
S=\frac32\pi.
\]
在圆心处
\[
x+y=1.
\]
故
\[
\iint_D(x+y)\,dxdy=1\cdot S=\frac{3\pi}{2}.
\]
\answer{\(\displaystyle \frac{3\pi}{2}\)}
\examnote{二次区域先配方；若被积函数是线性的，中心对称能直接给答案。}
\end{solutionblock}

\begin{problemblock}
\textbf{例6.} 设 \(f(x)\) 为连续奇函数，区域
\[
D=\{(x,y)\mid (x+1)^2+(y+1)^2\le1\},
\]
则（\quad）

A. \(\displaystyle \iint_D f(x+y)\,dxdy=0\)
\quad
B. \(\displaystyle \iint_D f(x-y)\,dxdy=0\)

C. \(\displaystyle \iint_D f(x^2y)\,dxdy=0\)
\quad
D. \(\displaystyle \iint_D f(xy^2)\,dxdy=0\)
\end{problemblock}
\begin{solutionblock}
\analysis{圆盘中心在 \((-1,-1)\)，但它关于直线 \(x=y\) 对称。只有 \(x-y\) 在交换 \(x,y\) 时变号。}
区域 \(D\) 满足：若 \((x,y)\in D\)，则 \((y,x)\in D\)。在变换
\[
(x,y)\mapsto(y,x)
\]
下，
\[
x-y\mapsto y-x=-(x-y).
\]
由于 \(f\) 是奇函数，
\[
f(y-x)=-f(x-y).
\]
区域和面积元不变，因此
\[
\iint_D f(x-y)\,dxdy
=-\iint_D f(x-y)\,dxdy,
\]
故该积分为 \(0\)。
\answer{B}
\examnote{判断“任意奇函数积分为零”，关键是找到保持区域不变且使函数自变量变号的变换。}
\end{solutionblock}

\section{换序、变量替换与证明}

\begin{problemblock}
\textbf{例7.} 计算
\[
\int_0^1dy\int_y^1
\left(\frac{e^y}{1+x^2}+\frac{e^x}{1+y^2}\right)dx.
\]
\begin{enumerate}[label=\Alph*., leftmargin=2.2em]
\item \(\displaystyle \frac{\pi(e-1)}4\)
\item \(\displaystyle \frac{\pi e}{4}\)
\item \(\displaystyle \frac{\pi(e-1)}2\)
\item \(\displaystyle \frac{\pi e}{2}\)
\end{enumerate}
\end{problemblock}
\begin{solutionblock}
\analysis{积分区域是三角形 \(0\le y\le x\le1\)。两个被积函数互为交换 \(x,y\) 后的对应项。}
记
\[
h(x,y)=\frac{e^y}{1+x^2}.
\]
原积分为
\[
\iint_{0\le y\le x\le1}\bigl[h(x,y)+h(y,x)\bigr]\,dxdy.
\]
由于两个三角形 \(0\le y\le x\le1\) 与 \(0\le x\le y\le1\) 互补成单位正方形，上式等于
\[
\iint_{[0,1]^2}\frac{e^y}{1+x^2}\,dxdy.
\]
分离变量：
\[
\int_0^1e^y\,dy\int_0^1\frac{dx}{1+x^2}
=(e-1)\cdot\frac\pi4.
\]
\answer{A}
\examnote{三角形上出现 \(h(x,y)+h(y,x)\)，可直接扩展为正方形上的积分。}
\end{solutionblock}

\begin{problemblock}
\textbf{例8.} 计算
\[
\iint_D\frac{\sin\sqrt{y^2-xy}}{y}\,dxdy,
\]
其中 \(D\) 是由直线 \(y=x\)、直线 \(y=1\) 和 \(y\) 轴围成的平面区域。
\end{problemblock}
\begin{solutionblock}
\analysis{区域为 \(0\le x\le y\le1\)。根号中 \(y^2-xy=y(y-x)\)，令 \(u=\sqrt{y^2-xy}\) 可化简。}
区域写为
\[
0\le y\le1,\qquad 0\le x\le y.
\]
令
\[
u=y,\qquad v=\sqrt{y^2-xy}.
\]
则
\[
y=u,\qquad x=u-\frac{v^2}{u}.
\]
当 \(0\le x\le y\le1\) 时，
\[
0\le v\le u\le1.
\]
Jacobian 为
\[
\left|\frac{\partial(x,y)}{\partial(u,v)}\right|
=\frac{2v}{u}.
\]
原积分变为
\[
\int_0^1\int_0^u \frac{\sin v}{u}\cdot\frac{2v}{u}\,dvdu.
\]
换序：
\[
\int_0^1 2v\sin v\int_v^1\frac{du}{u^2}\,dv
=\int_0^1 2v\sin v\left(\frac1v-1\right)dv.
\]
即
\[
2\int_0^1(1-v)\sin v\,dv
=2(1-\sin1).
\]
\answer{\(\displaystyle 2(1-\sin1)\)}
\examnote{根式 \(\sqrt{y(y-x)}\) 与区域 \(0\le x\le y\) 匹配时，可令新变量为该根式。}
\end{solutionblock}

\begin{problemblock}
\textbf{例9.} 设
\[
f(t)=\int_t^{2t}dx\int_x^t e^{(x-y+1)^2}\,dy,
\]
则
\[
\lim_{t\to0^+}\frac{f(t)}{t^2}=(\quad)
\]

A. \(\dfrac e2\)\qquad B. \(-\dfrac e2\)\qquad C. \(2e\)\qquad D. \(-2e\)
\end{problemblock}
\begin{solutionblock}
\analysis{积分区域随 \(t\) 缩小，令 \(x=ta,\ y=tb\)，区域固定后看被积函数极限。注意内层积分上限 \(t\) 小于 \(x\)，带负向。}
令
\[
x=ta,\qquad y=tb.
\]
则 \(a\in[1,2]\)，内层 \(b\) 从 \(a\) 到 \(1\)，且 \(dxdy=t^2dadb\)。于是
\[
\frac{f(t)}{t^2}
=\int_1^2 da\int_a^1 e^{(t(a-b)+1)^2}\,db.
\]
令 \(t\to0^+\)，被积函数趋于 \(e\)，故
\[
\lim_{t\to0^+}\frac{f(t)}{t^2}
=e\int_1^2(1-a)\,da
=-\frac e2.
\]
\answer{B}
\examnote{变区域极限常用伸缩变换；积分方向若反向，符号最容易错。}
\end{solutionblock}

\begin{problemblock}
\textbf{例10.} 设
\[
D=\{(x,y)\mid 0\le x\le2,\ 0\le y\le2\}.
\]
\begin{enumerate}[label=(\arabic*), leftmargin=2.2em]
\item 求
\[
A=\iint_D|xy-1|\,dxdy;
\]
\item 设 \(f(x,y)\) 在 \(D\) 上连续，且
\[
\iint_Df(x,y)\,dxdy=0,\qquad
\iint_Dxyf(x,y)\,dxdy=1.
\]
证明：存在 \((\xi,\eta)\in D\)，使
\[
|f(\xi,\eta)|\ge\frac1A.
\]
\end{enumerate}
\end{problemblock}
\begin{solutionblock}
\analysis{第 (1) 问按曲线 \(xy=1\) 分区；第 (2) 问把两个积分条件相减，得到 \(\iint_D(xy-1)f=1\)。}
(1) 当 \(0\le x\le1/2\) 时，\(0\le xy\le1\)，内层积分为
\[
\int_0^2(1-xy)\,dy=2-2x.
\]
当 \(1/2\le x\le2\) 时，以 \(y=1/x\) 分段：
\[
\int_0^2|xy-1|\,dy
=\int_0^{1/x}(1-xy)\,dy+\int_{1/x}^{2}(xy-1)\,dy
=2x-2+\frac1x.
\]
所以
\[
\begin{aligned}
A&=\int_0^{1/2}(2-2x)\,dx
+\int_{1/2}^{2}\left(2x-2+\frac1x\right)dx\\
&=\frac34+\frac34+2\ln2
=\frac32+2\ln2.
\end{aligned}
\]

(2) 由题设
\[
1=\iint_Dxyf(x,y)\,dxdy-\iint_Df(x,y)\,dxdy
=\iint_D(xy-1)f(x,y)\,dxdy.
\]
于是
\[
1\le \iint_D|xy-1|\,|f(x,y)|\,dxdy
\le \max_D|f(x,y)|\iint_D|xy-1|\,dxdy.
\]
即
\[
\max_D|f(x,y)|\ge\frac1A.
\]
因 \(f\) 连续，\(|f|\) 在闭区域 \(D\) 上取到最大值，故存在 \((\xi,\eta)\in D\)，使
\[
|f(\xi,\eta)|\ge\frac1A.
\]
\answer{(1) \(\displaystyle A=\frac32+2\ln2\)；(2) 已证。}
\examnote{证明“存在一点使函数值不小于某数”时，常用积分估计反推最大值下界。}
\end{solutionblock}

\begin{problemblock}
\textbf{例11.} 设 \(f(x,y)\) 为连续函数，求
\[
\int_1^2dx\int_0^{\sqrt{2x-x^2}}f(x,y)\,dy
\]
的等价积分表示。
\end{problemblock}
\begin{solutionblock}
\analysis{区域为圆 \((x-1)^2+y^2\le1\) 的上半部分中 \(x\ge1\) 的右半块。极坐标下边界为 \(r=2\cos\theta\)，并且 \(x\ge1\) 即 \(r\cos\theta\ge1\)。}
区域满足
\[
(x-1)^2+y^2\le1,\qquad y\ge0,\qquad x\ge1.
\]
用极坐标
\[
x=r\cos\theta,\qquad y=r\sin\theta.
\]
圆边界为
\[
r^2=2r\cos\theta,\qquad r=2\cos\theta.
\]
条件 \(x\ge1\) 为
\[
r\cos\theta\ge1,\qquad r\ge\sec\theta.
\]
又 \(y\ge0\)，且交点方向给出
\[
0\le\theta\le\frac\pi4.
\]
所以
\[
\int_1^2dx\int_0^{\sqrt{2x-x^2}}f(x,y)\,dy
=\int_0^{\pi/4}d\theta
\int_{\sec\theta}^{2\cos\theta}
f(r\cos\theta,r\sin\theta)r\,dr.
\]
\answer{\(\displaystyle \int_0^{\pi/4}\int_{\sec\theta}^{2\cos\theta} f(r\cos\theta,r\sin\theta)r\,dr\,d\theta\)}
\examnote{极坐标中不要漏掉 \(x\ge1\) 产生的下界 \(r=\sec\theta\)。}
\end{solutionblock}

\section{特殊变量替换与复杂区域}

\begin{problemblock}
\textbf{例12.} 设平面有界区域 \(D\) 位于第一象限，由
\[
xy=\frac13,\quad xy=3,\quad y=\frac13x,\quad y=3x
\]
围成。计算
\[
\iint_D(1+x-y)\,dxdy.
\]
\end{problemblock}
\begin{solutionblock}
\analysis{边界正好提示令 \(u=xy,\ v=y/x\)。}
令
\[
u=xy,\qquad v=\frac yx.
\]
在第一象限中
\[
x=\sqrt{\frac uv},\qquad y=\sqrt{uv}.
\]
边界变为
\[
\frac13\le u\le3,\qquad \frac13\le v\le3.
\]
Jacobian：
\[
\frac{\partial(u,v)}{\partial(x,y)}
=2v,
\qquad
dxdy=\frac1{2v}\,dudv.
\]
于是
\[
\iint_D(1+x-y)\,dxdy
=\int_{1/3}^3\int_{1/3}^3
\left(1+\sqrt{\frac uv}-\sqrt{uv}\right)\frac1{2v}\,dvdu.
\]
其中
\[
\int_{1/3}^3v^{-3/2}\,dv
=\int_{1/3}^3v^{-1/2}\,dv=\frac4{\sqrt3},
\]
所以含 \(\sqrt u\) 的两项相消。剩下
\[
\frac12\int_{1/3}^3du\int_{1/3}^3\frac{dv}{v}
=\frac12\cdot\frac83\cdot\ln9
=\frac83\ln3.
\]
\answer{\(\displaystyle \frac83\ln3\)}
\examnote{双曲线 \(xy=\text{常数}\) 与射线 \(y=kx\) 围成区域，标准换元是 \(u=xy,\ v=y/x\)。}
\end{solutionblock}

\begin{problemblock}
\textbf{例13.}（仅数一、数二）设平面有界区域 \(D\) 位于第一象限，由曲线
\[
x^2+y^2-xy=1,\qquad x^2+y^2-xy=2
\]
与直线
\[
y=\sqrt3x,\qquad y=0
\]
围成。计算
\[
\iint_D\frac1{3x^2+y^2}\,dxdy.
\]
\end{problemblock}
\begin{solutionblock}
\analysis{曲线含 \(x^2+y^2-xy\)，在极坐标下变为 \(r^2(1-\frac12\sin2\theta)\)。}
令
\[
x=r\cos\theta,\qquad y=r\sin\theta.
\]
区域角度为
\[
0\le\theta\le\frac\pi3.
\]
且
\[
x^2+y^2-xy=r^2(1-\sin\theta\cos\theta).
\]
对固定 \(\theta\)，\(r\) 从
\[
\sqrt{\frac1{1-\sin\theta\cos\theta}}
\]
到
\[
\sqrt{\frac2{1-\sin\theta\cos\theta}}.
\]
又
\[
3x^2+y^2=r^2(3\cos^2\theta+\sin^2\theta),
\]
所以
\[
\begin{aligned}
I&=\int_0^{\pi/3}
\frac{1}{3\cos^2\theta+\sin^2\theta}
\int_{r_1}^{r_2}\frac{dr}{r}\,d\theta\\
&=\frac12\ln2\int_0^{\pi/3}
\frac{d\theta}{3\cos^2\theta+\sin^2\theta}.
\end{aligned}
\]
令 \(t=\tan\theta\)，则
\[
\int_0^{\pi/3}
\frac{d\theta}{3\cos^2\theta+\sin^2\theta}
=\int_0^{\sqrt3}\frac{dt}{t^2+3}
=\frac{\pi}{4\sqrt3}.
\]
因此
\[
I=\frac{\pi\ln2}{8\sqrt3}.
\]
\answer{\(\displaystyle \frac{\pi\ln2}{8\sqrt3}\)}
\examnote{二次型区域常用极坐标；若只含角度的因子，径向积分往往给出 \(\ln\) 常数。}
\end{solutionblock}

\begin{problemblock}
\textbf{例14.} 设有界区域 \(D\) 是由圆 \(x^2+y^2=1\)、直线 \(y=x\) 以及 \(x\) 轴在第一象限内围成的部分，计算
\[
\iint_D e^{(x+y)^2}(x^2-y^2)\,dxdy.
\]
\end{problemblock}
\begin{solutionblock}
\analysis{区域是扇形 \(0\le\theta\le\pi/4,\ 0\le r\le1\)。被积函数中的 \(x^2-y^2\) 与 \((x+y)^2\) 在极坐标下可配合换元。}
令
\[
x=r\cos\theta,\qquad y=r\sin\theta.
\]
则
\[
0\le r\le1,\qquad 0\le\theta\le\frac\pi4,
\]
\[
x^2-y^2=r^2\cos2\theta,
\]
\[
(x+y)^2=r^2(\cos\theta+\sin\theta)^2=r^2(1+\sin2\theta).
\]
原积分为
\[
\int_0^1r^3dr\int_0^{\pi/4}
e^{r^2(1+\sin2\theta)}\cos2\theta\,d\theta.
\]
令
\[
w=1+\sin2\theta,\qquad dw=2\cos2\theta\,d\theta.
\]
当 \(\theta=0\) 时 \(w=1\)，当 \(\theta=\pi/4\) 时 \(w=2\)。故内层积分为
\[
\frac12\int_1^2e^{r^2w}\,dw
=\frac{e^{2r^2}-e^{r^2}}{2r^2}.
\]
于是原积分
\[
\frac12\int_0^1r(e^{2r^2}-e^{r^2})\,dr
=\frac18(e-1)^2.
\]
\answer{\(\displaystyle \frac{(e-1)^2}{8}\)}
\examnote{看到 \(x^2-y^2\) 和 \((x+y)^2\)，极坐标后用 \(1+\sin2\theta\) 换元非常顺。}
\end{solutionblock}

\begin{problemblock}
\textbf{例15.} 设平面区域 \(D\) 由曲线
\[
x=t-\sin t,\qquad y=1-\cos t,\qquad 0\le t\le2\pi
\]
与 \(x\) 轴围成。计算
\[
\iint_D(x+2y)\,d\sigma.
\]
\end{problemblock}
\begin{solutionblock}
\analysis{曲线是一个摆线拱。因 \(dx/dt=1-\cos t=y\ge0\)，可用竖条积分。}
记
\[
X(t)=t-\sin t,\qquad Y(t)=1-\cos t.
\]
区域可按 \(0\le y\le Y(t)\) 竖向积分，且
\[
dx=Y(t)\,dt.
\]
对固定 \(x=X(t)\)，
\[
\int_0^{Y(t)}(x+2y)\,dy
=X(t)Y(t)+Y^2(t).
\]
故
\[
I=\int_0^{2\pi}\bigl[X(t)Y(t)+Y^2(t)\bigr]Y(t)\,dt.
\]
即
\[
I=\int_0^{2\pi}(t-\sin t)(1-\cos t)^2\,dt
+\int_0^{2\pi}(1-\cos t)^3\,dt.
\]
利用周期对称性，
\[
\int_0^{2\pi}t(1-\cos t)^2dt
=\pi\int_0^{2\pi}(1-\cos t)^2dt=3\pi^2,
\]
且
\[
\int_0^{2\pi}\sin t(1-\cos t)^2dt=0,
\]
\[
\int_0^{2\pi}(1-\cos t)^3dt=5\pi.
\]
因此
\[
I=3\pi^2+5\pi.
\]
\answer{\(\displaystyle 3\pi^2+5\pi\)}
\examnote{参数曲线围成区域可用 \(dx=x'(t)dt\) 的竖条法；摆线中 \(x'(t)=y(t)\) 是关键。}
\end{solutionblock}

\begin{problemblock}
\textbf{例16.} 设二元函数
\[
f(x,y)=
\begin{cases}
x^2,& |x|+|y|\le1,\\[2mm]
\dfrac1{\sqrt{x^2+y^2}},&1<|x|+|y|\le2.
\end{cases}
\]
计算
\[
\iint_D f(x,y)\,d\sigma,\qquad
D=\{(x,y)\mid |x|+|y|\le2\}.
\]
\end{problemblock}
\begin{solutionblock}
\analysis{区域是菱形。内层菱形直接算 \(x^2\)，外层用极坐标下的菱形边界 \(r(|\cos\theta|+|\sin\theta|)=a\)。}
内层
\[
D_1=\{|x|+|y|\le1\}
\]
上
\[
\iint_{D_1}x^2\,d\sigma
=4\int_0^1\int_0^{1-x}x^2\,dy\,dx
=4\int_0^1x^2(1-x)\,dx=\frac13.
\]
外层为 \(\{|x|+|y|\le2\}\setminus\{|x|+|y|\le1\}\)。在第一象限中，边界为
\[
r(\cos\theta+\sin\theta)=a.
\]
由于被积函数为 \(1/r\)，面积元为 \(rdrd\theta\)，外层积分为
\[
4\int_0^{\pi/2}\int_{\frac1{\cos\theta+\sin\theta}}^{\frac2{\cos\theta+\sin\theta}}dr\,d\theta.
\]
即
\[
4\int_0^{\pi/2}\frac{d\theta}{\cos\theta+\sin\theta}.
\]
而
\[
\cos\theta+\sin\theta=\sqrt2\cos\left(\theta-\frac\pi4\right),
\]
所以外层积分为
\[
4\sqrt2\ln(1+\sqrt2).
\]
总积分为
\[
\frac13+4\sqrt2\ln(1+\sqrt2).
\]
\answer{\(\displaystyle \frac13+4\sqrt2\ln(1+\sqrt2)\)}
\examnote{\(|x|+|y|\le a\) 的菱形在极坐标下边界是 \(r=a/(|\cos\theta|+|\sin\theta|)\)。}
\end{solutionblock}

\begin{problemblock}
\textbf{例17.} 计算
\[
I=\iint_D |y-x^2|\max\{x,y\}\,d\sigma,
\qquad
D=\{(x,y)\mid 0\le x\le1,\ 0\le y\le1\}.
\]
\end{problemblock}
\begin{solutionblock}
\analysis{在单位正方形内有 \(x^2\le x\)。按 \(y=x^2\) 和 \(y=x\) 分成三块。}
对固定 \(x\in[0,1]\)，分三段：
\[
0\le y\le x^2,\qquad x^2\le y\le x,\qquad x\le y\le1.
\]
于是
\[
\begin{aligned}
I=\int_0^1&\left[
\int_0^{x^2}(x^2-y)x\,dy
+\int_{x^2}^{x}(y-x^2)x\,dy\right.\\
&\left.+\int_x^1(y-x^2)y\,dy
\right]dx.
\end{aligned}
\]
逐项计算后，内层总和为
\[
\frac13-\frac{x^2}{2}+\frac{x^3}{6}-\frac{x^4}{2}+x^5.
\]
因此
\[
I=\int_0^1\left(
\frac13-\frac{x^2}{2}+\frac{x^3}{6}-\frac{x^4}{2}+x^5
\right)dx
=\frac{11}{40}.
\]
\answer{\(\displaystyle \frac{11}{40}\)}
\examnote{含绝对值和 \(\max\) 的二重积分，先画出分界曲线；本题分界为 \(y=x^2\) 与 \(y=x\)。}
\end{solutionblock}

\begin{problemblock}
\textbf{例18.} 已知平面区域
\[
D=\{(x,y)\mid (x-1)^2+y^2\le1\},
\]
计算
\[
\iint_D\left|\sqrt{x^2+y^2}-1\right|\,dxdy.
\]
\end{problemblock}
\begin{solutionblock}
\analysis{区域是以 \((1,0)\) 为圆心、半径为 \(1\) 的圆盘。用以原点为极点的极坐标，边界为 \(r=2\cos\theta\)。}
在极坐标中，
\[
D:\quad -\frac\pi2\le\theta\le\frac\pi2,\qquad 0\le r\le2\cos\theta.
\]
积分为
\[
I=\int_{-\pi/2}^{\pi/2}\int_0^{2\cos\theta}|r-1|r\,drd\theta.
\]
当 \(|\theta|\le\pi/3\) 时，\(2\cos\theta\ge1\)，需在 \(r=1\) 分段；当 \(\pi/3<|\theta|\le\pi/2\) 时，\(2\cos\theta<1\)。利用偶对称：
\[
\begin{aligned}
I=2\bigg[&
\int_0^{\pi/3}\left(\int_0^1(1-r)r\,dr+\int_1^{2\cos\theta}(r-1)r\,dr\right)d\theta\\
&+\int_{\pi/3}^{\pi/2}\int_0^{2\cos\theta}(1-r)r\,drd\theta
\bigg].
\end{aligned}
\]
计算得
\[
I=3\sqrt3-\frac{\pi}{9}-\frac{32}{9}.
\]
\answer{\(\displaystyle 3\sqrt3-\frac{\pi}{9}-\frac{32}{9}\)}
\examnote{偏心圆盘遇到 \(\sqrt{x^2+y^2}\) 时，以原点为极点更自然，但要按 \(r=1\) 分段。}
\end{solutionblock}

\begin{problemblock}
\textbf{例19.} 设闭区域
\[
D=\{(x,y)\mid0\le x\le1,\ 0\le y\le1\},
\]
\(f(x,y)\) 为 \(D\) 上的连续函数，且
\[
f(x,y)=\sqrt{x^2+y^2}\max\{x,y\}
+2\int_0^1f(x,0)\,dx+\iint_Df(x,y)\,dxdy.
\]
求 \(f(x,y)\)。
\end{problemblock}
\begin{solutionblock}
\analysis{右侧后两项是常数。设其为 \(C\)，把问题化成常数方程。}
记
\[
h(x,y)=\sqrt{x^2+y^2}\max\{x,y\},
\]
并设
\[
C=2\int_0^1f(x,0)\,dx+\iint_Df(x,y)\,dxdy.
\]
则
\[
f(x,y)=h(x,y)+C.
\]
先计算
\[
\int_0^1 f(x,0)\,dx
=\int_0^1 x^2\,dx+C=\frac13+C.
\]
再记
\[
H=\iint_D h(x,y)\,dxdy.
\]
利用对称性，
\[
H=2\int_0^1\int_0^x x\sqrt{x^2+y^2}\,dy\,dx.
\]
令 \(y=xt\)，得
\[
H=2\int_0^1x^3dx\int_0^1\sqrt{1+t^2}\,dt
=\frac14\left[\sqrt2+\ln(1+\sqrt2)\right].
\]
因此
\[
\iint_Df=H+C.
\]
代回 \(C\) 的定义：
\[
C=2\left(\frac13+C\right)+(H+C).
\]
故
\[
-2C=\frac23+H,
\qquad
C=-\frac13-\frac H2.
\]
即
\[
C=-\frac13-\frac18\left[\sqrt2+\ln(1+\sqrt2)\right].
\]
所以
\[
f(x,y)=\sqrt{x^2+y^2}\max\{x,y\}
-\frac13-\frac18\left[\sqrt2+\ln(1+\sqrt2)\right].
\]
\answer{\(\displaystyle f(x,y)=\sqrt{x^2+y^2}\max\{x,y\}
-\frac13-\frac18\bigl[\sqrt2+\ln(1+\sqrt2)\bigr]\)}
\examnote{含未知函数积分的方程，先把积分项整体设为常数，再反代求常数。}
\end{solutionblock}
